\section{Relative permanent estimates with a growing entry bound}
\label{sec:relativepermanents}
The preceding entropy bounds do not determine a relative
\(1+o(1)\) count.
For that purpose the entry parameter in a permanent asymptotic
must be tracked; a theorem uniform only for a fixed bound
cannot simply be applied when that bound grows.
Throughout this section \(J_N\) is the all-ones matrix.

\begin{theorem}[Polylogarithmic entry bounds in the permanent expansion]
\label{thm:growingpermanent}
Fix \(\delta\in(0,1)\) and \(K>0\).
Let \(D\) be real nonnegative doubly stochastic of order \(N\),
and write
\[
 B=D-J_N/N,\qquad
 \|B\|\leq1-\delta,\qquad
 \max_{ij}D_{ij}\leq A/N,\qquad 2\leq A\leq(\log N)^K .
\]
Then, uniformly over these matrices for sufficiently large \(N\),
\[
 \operatorname{per}D
 =\frac{N!}{N^N}\det(I-B^\top B)^{-1/2}
                       \bigl(1+O_{\delta,K}(N^{-1/8})\bigr).
\]
Zero entries are allowed.
\end{theorem}

\begin{proof}
We track the finite estimates behind the fixed-entry theorem
of Divoux--Wang~\cite[Section 3]{DW}, based on McCullagh~\cite{McC}.
This is a parameter-uniformity argument, not an application
of their fixed-\(A\) asymptotic statement.

Stochasticity gives the useful stronger bounds
\[
 \max_i\sum_jB_{ij}^2,\quad
 \max_j\sum_iB_{ij}^2\leq A/N,\qquad
 \|B\|_F^2\leq A.
\]
Choose fixed \(0<a<1<b\) and \(R>1\) so that
\(bR(1-\delta/2)<1\).
Let \(R_B(s)\) be the rook polynomial
\[
 R_B(s)=\sum_{k=0}^N s^k
       \sum_{\substack{I,J\subseteq[1,N]\\|I|=|J|=k}}
                         \operatorname{per}B[I,J].
\]
The exact rook/Gamma identity, with
\(T\sim\Gamma(N+1,1)\), is
\[
 \frac{N^N\operatorname{per}D}{N!}
                         =\mathbb E R_B(N/T).
\]
Define
\[
 F_s(t)=R_{tB}(s)
             \det(I-s^2t^2B^\top B)^{1/2},
 \qquad s\in[a,b],
\]
with the analytic square root equal to one at \(t=0\).
We show that \(F_s(1)=1+O_\delta(A^2N^{-1/6})\),
uniformly in \(s\).

Here are the finite Gaussian estimates and their parameter
dependence explicitly.
With
\[
 \Omega_{s,t}=\begin{pmatrix}0&stB\\stB^\top&0\end{pmatrix},
 \quad
 \Gamma_{s,t}=\Omega_{s,t}(I+\Omega_{s,t})^{-1},
\]
the normalized Gaussian identity in~\cite[Lemmas 3.5--3.7]{DW}
writes \(F_s(t)\) as the Gaussian expectation of
\(\prod_i r(Y_i)\), where
\[
 r(z)=(1+z)\exp(-z+z^2/2)=1+O(z^3)
\]
and \(Y\) has bilinear covariance \(\Gamma_{s,t}\).
The zero row and column sums of \(B\) cancel the linear
Gaussian term exactly, before taking any bound.

The polynomial transfer underlying this identity has an independent
maintained proof. For any finite coordinate set, choose complex numbers
\(z_j,w_j\) with \(|z_j|<1\) and
\(w_j^2=\gamma_j=z_j/(1+z_j)\), and put
\(b_j=(1-\gamma_j)/2\). The complete quotient gives
\[
 (1-2\Re\gamma_j)|1+z_j|^2=1-|z_j|^2,
 \qquad \Re b_j>1/4.
\]
Every polynomial times the resulting complex Gaussian density is
integrable, and each scalar normalizer is nonzero. Define the actual
normalized integral
\[
 \mathcal G_z(P)=
 \frac{\displaystyle\int_{\mathbb R^J}
      P((w_jx_j)_j)\exp(-\sum_jb_jx_j^2)\,dx}
      {\displaystyle\prod_j\int_{\mathbb R}\exp(-b_jx^2)\,dx}.
\]
Scalar integration by parts and finite-dimensional Fubini prove
\[
 \mathcal G_z(X^d)=\prod_j W_{z_j}(d_j),\quad
 W_z(2e)=(2e-1)!!z^e,\quad W_z(2e+1)=0,\quad W_z(0)=1,
\]
and, for every polynomial \(P\),
\(\mathcal G_z(X_jP)=z_j\mathcal G_z(\partial_jP)\).
For \(L_a=\sum_ja_jX_j\), set
\(C_z(a,b)=\sum_jz_ja_jb_j\). The complete recursion is
\[
 \mathcal G_z\left(L_a\prod_{\ell=1}^mL_{a_\ell}\right)
 =\sum_{k=1}^m C_z(a,a_k)\,
    \mathcal G_z\left(\prod_{\substack{1\leq\ell\leq m\\\ell\ne k}}
                                      L_{a_\ell}\right).
\]
In particular, if every covariance among \(2e\) listed forms has
modulus at most \(c\geq0\), the actual normalized integral satisfies
\[
 \left|\mathcal G_z\left(\prod_{\ell=1}^{2e}L_{a_\ell}\right)\right|
                         \leq(2e-1)!!c^e .
\]
All these statements, including integrability, all repeated-vertex
contractions and the uniform bound in \(e\), have maintained Lean
proofs. The bound is used on the diagram remainder after the signed
low-degree terms have been retained.

The restriction \(|z_j|<1\) belongs to the quadratic transfer, not
to the existence of a Gaussian realization. For arbitrary complex
\(v_j\), choose \(u_j^2=v_j\). Standard Gaussian integration gives
\[
 \frac{\displaystyle\int_{\mathbb R^J}
      P((u_jx_j)_j)e^{-\sum_jx_j^2/2}\,dx}
      {Z_0^{|J|}}=\mathcal G_v(P),\qquad
 Z_0=\int_{\mathbb R}e^{-x^2/2}\,dx=\sqrt{2\pi}.
\]
This unrestricted realization, including the equality of the real and
complex scalar normalizations, is formalized. Thus an indefinite
or large compensated covariance is not excluded by a spurious
positive-semidefinite or norm restriction.

The exact rook identification is also formalized. Let \(L_i,R_j\),
\(1\leq i,j\leq N\), be linear forms whose within-left and within-right
covariances vanish. For every \(S\subseteq\{1,\ldots,N\}\), the maintained
support identity is
\[
 (N-|S|)!\,\mathcal G_v\!\left(
       \prod_{i\in S}L_i\prod_{j=1}^{N}(1+R_j)\right)
 =\sum_{\sigma\in\mathfrak S_N}
       \prod_{i\in S}\mathcal G_v(L_iR_{\sigma(i)}).
\]
For nonempty \(S\), delete one selected vertex \(i\) and contract it
with each right vertex \(j\). The support loses one element and
\(N\) decreases by one, so the factorial is identical on every branch.
The full permutation decomposition at column \(i\) has exactly these
branches. For empty \(S\), the Gaussian expectation is one and
the permutation count is \(N!\). Summing every support and dividing
by its nonzero factorial gives the complete signed affine-product
rook formula, without a discarded contraction or an assumed independence.

For a real symmetric matrix \(M=U\diag(\lambda_j)U^\top\), let
\(Y_i=\sum_jU_{ij}\sqrt{\lambda_j}X_j\), where the square roots
are complex and the \(X_j\) are independent real standard normals.
The maintained spectral proof derives
\[
 \mathcal G_\lambda(L_iL_k)=M_{ik},\qquad
 W:=\sum_i|Y_i|^2=\sum_j|\lambda_j|X_j^2.
\]
If \(\one^\top M=0\), it also proves \(\sum_iY_i=0\) pointwise.
Both identities retain every cross term until orthogonality cancels it.
For the actual \(M=\Omega\) from a doubly stochastic \(D\),
symmetry and this zero-sum condition are discharged directly.

The scalar Gaussian integral and finite-dimensional Fubini give, for
\(\tau\geq0\), \(\|M\|\leq\rho\) and \(\tau\rho<1\),
\[
 \E e^{\tau W/2}
 \leq\exp\!\left(\frac{\tau}{1-\tau\rho}\sum_j|\lambda_j|\right).
\]
The exact Frobenius identity
\(\sum_j\lambda_j^2=\|M\|_F^2\) is formalized.
For \(M=\Omega\), stochasticity gives
\(\|M\|_F^2=2\|B\|_F^2\leq2(A-1)\), hence
\[
 \E e^{\tau W/2}
 \leq\exp\!\left(\frac{2\tau}{1-\tau\rho}
                                      \sqrt{N(A-1)}\right).
\]
Finally the exact zero-sum product bound, with \(c=\sqrt t\),
proves the actual collective integral estimate
\[
 \boxed{\quad
 \E\prod_{i=1}^{2N}|1+\sqrt t\,Y_i|
 \leq\exp\!\left(\frac{2|t|}{1-|t|\rho}
                                     \sqrt{N(A-1)}\right),
 \quad |t|\rho<1.\quad}
\]
Every hypothesis in this stochastic consumer, including integrability,
is discharged in maintained Lean. This is the Cauchy/tail input;
the signed low-degree corrections are retained separately.
The exact Gaussian-to-rook consumer now transfers this bound directly:
\[
 \boxed{\quad |R_B(t)|\leq
 \exp\!\left(\frac{2|t|}{1-|t|\rho}\sqrt{N(A-1)}\right),
 \qquad |t|\rho<1.\quad}
\]
The maintained complex polynomial agrees with the original real rook
polynomial at every real parameter. The compensated resolvent
identification and the complete diagram, Cauchy and outside-Gamma
consumers remain formalization obligations.

The block-power expansion gives, for a fixed
\(\theta=b(1-\delta/2)<1\),
\[
 \max_{i,j}|[t^\ell](\Gamma_{s,t})_{ij}|
       \leq(2b^2 A/N)\theta^{\ell-1},
 \qquad \ell\geq1.
\]
Indeed each power of length at least two is bounded using
the two row/column Euclidean norms above and the intervening
operator power. The length-one term is at most \(bA/N\).
Thus the covariance constant grows linearly in \(A\),
not as an unspecified function. The actual real matrix geometric
series and this bound for every degree are formalized; the complex
analytic identification used here is a further consumer.

For clarity the finite Wick coefficient estimate used here is:
if \(F_s(t)=1+\sum_{k\geq1}f_k(s)t^k\), then
\[
 |f_k(s)|\leq
 \sum_{e=1}^k\ \sum_{1\leq v\leq\lfloor2e/3\rfloor}
 (C_\delta A)^e(e/N)^{e-v}
                  \binom{k-1}{e-1}\theta^{k-e}.
\]
This is the explicit estimate in~\cite[Lemma 3.7]{DW}.
Its constant is an absolute constant times the covariance
constant: Wick pairings have \(e\) pairs, the nonconstant
vertices have degree at least three, hence \(v\leq2e/3\),
and their choices supply \(N^v\) against \(N^{-e}\).
There is no hidden fixed-entry asymptotic in this finite bound.

Take \(H=\lceil C_\delta A^3\sqrt N\rceil\), increasing its
constant below if necessary. For large \(N\), \(H<N\).
Sum over \(k\leq H\), use
\(\sum_{k\geq e}\binom{k-1}{e-1}\theta^{k-e}
=(1-\theta)^{-e}\), and retain every degree choice.
Since \(e/N\leq H/N\) and \(e-v\geq e/3\),
\[
 \sum_{1\leq k\leq H}|f_k(s)|
 \leq\sum_{e\geq1}e
       \bigl(C_\delta A^2N^{-1/6}\bigr)^e
 =O_\delta(A^2N^{-1/6}).
\]
The base tends to zero throughout the prescribed range of \(A\).

The same Gaussian product estimate on \(|t|=R\) gives
\[
 \max_{|t|=R}|F_s(t)|
 \leq\exp\bigl(C_\delta\sqrt{NA}+C_\delta A\bigr).
\]
To see its dependence, the Gaussian exponential moment is
bounded by a product \(\prod_\ell(1-bR\sigma_\ell)^{-1}\);
its logarithm is at most
\(C_\delta\sum_\ell\sigma_\ell\leq C_\delta\sqrt{NA}\).
The additional determinant factor costs at most
\(C_\delta\sum_\ell\sigma_\ell^2\leq C_\delta A\).
Cauchy's coefficient bound and the chosen \(H\) now give
\[
 \sum_{k>H}|f_k(s)|
       \leq\exp(-c_\delta A^3\sqrt N).
\]
Consequently the asserted relative rook estimate follows.

It remains to average the determinant factor over \(T\);
replacing it by an absolute bounded constant would lose
uniformity. Put
\(\Delta(s)=\det(I-s^2B^\top B)^{-1/2}\).
On \([a,b]\),
\[
 0\leq\log\Delta(s)\leq C_\delta A,\qquad
 |(\log\Delta)'(s)|\leq C_\delta A .
\]
On \(|N/T-1|\leq N^{-1/4}\), this gives
\(\Delta(N/T)/\Delta(1)=1+O_\delta(AN^{-1/4})\).
The remaining Gamma mass inside \([a,b]\) is at most
\(\exp(-c_\delta\sqrt N)\), while the determinant ratio is
at most \(\exp(C_\delta A)\).

For completeness the outside-\([a,b]\) Gaussian/Gamma tail
has absolute bound
\[
 \mathbb E\bigl[\one_{\{N/T\notin[a,b]\}}
                         |R_B(N/T)|\bigr]
 \leq\exp(-c_\delta N+C_\delta\sqrt{NA}).
\]
Here is the collective cancellation underlying the change of
variable. In a real standard-Gaussian representation let
\(Y\in\mathbb C^{2N}\), \(\sum_iY_i=0\), and
\(W=\sum_i|Y_i|^2\). For every complex \(c\),
\[
 \sum_i|1+cY_i|^2=2N+|c|^2W,\qquad
 \prod_i|1+cY_i|\leq\exp(|c|^2W/2).
\]
The linear term vanishes before either bound is taken. Arithmetic
and geometric means, with \(c=\sqrt{N/t}\), also give
\[
 \prod_i|1+\sqrt{N/t}Y_i|
       \leq(1+W/(2t))^N\qquad(t>0).
\]
Multiplication by the actual Gamma density cancels its full power:
\[
 \frac{e^{-t}t^N}{N!}\prod_i|1+\sqrt{N/t}Y_i|
 \leq\frac{e^{W/2}e^{-(t+W/2)}(t+W/2)^N}{N!}.
\]
These deterministic inequalities have maintained Lean proofs.
Set \(w=W/2\) and shift \(u=t+w\) in the remaining integral.
On \(w\leq\eta N\), a sufficiently small fixed \(\eta>0\)
leaves a fixed Gamma-tail separation. On \(w>\eta N\), use
\[
 \E[e^w\one_{\{w>\eta N\}}]
 \leq e^{-(\lambda-1)\eta N}\E e^{\lambda w},
 \qquad 1<\lambda<(1-\delta)^{-1}.
\]
Diagonalizing the actual block matrix gives
\(W=\sum_\ell|\lambda_\ell|Z_\ell^2\) for independent
real standard normals. The paired eigenvalues yield
\(\E e^{\lambda W/2}
 =\prod_\ell(1-\lambda\sigma_\ell)^{-1}
 \leq\exp(C_\delta\sum_\ell\sigma_\ell)\).
This proves the stated tail with its Frobenius dependence retained,
as in~\cite[Section 3.3]{DW}. The Gaussian representation,
integration and tail consumer are still formalization obligations.
They vanish uniformly for polylogarithmic \(A\).

As \(\Delta(1)\geq1\), the total relative error is bounded by
\[
 C_\delta(A^2N^{-1/6}+AN^{-1/4})
 +\exp(C_\delta A-c_\delta\sqrt N)
 +\exp(-c_\delta N+C_\delta\sqrt{NA}).
\]
For fixed \(K\) this is \(O_{\delta,K}(N^{-1/8})\), as claimed.
\end{proof}

\begin{lemma}[Uniform compensated minors with mismatched supports]
\label{lem:relativecompensatedminors}
Under Theorem~\ref{thm:growingpermanent}, uniformly for arbitrary
row and column sets \(I,J\) with
\(|I|=|J|=r\leq(\log(2N))^3\),
\[
 \frac{\operatorname{per}D[I^c,J^c]}{\operatorname{per}D}
                  =1+O_{\delta,K}(N^{-1/9}).
\]
The removed supports need not coincide.
\end{lemma}
\begin{proof}
The case \(r=0\) is immediate. Otherwise put
\(n=N-r\), \(Q=D[I^c,J^c]\), and
\(s_{I,J}=\sum_{i\in I,j\in J}D_{ij}\).
Each remaining row or column loses at most
\(\varepsilon=rA/N\), and its total lost mass is exactly
\(r-s_{I,J}\). Also \(s_{I,J}\leq r^2A/N\).

The second singular value of \(Q\), after degree normalization,
stays a fixed distance from 1.
Indeed restrict \(B\) and compare \(J_n/N\) with \(J_n/n\);
their rank-one difference has norm \(r/N\).
The row and column normalization changes the operator by
\(O(\varepsilon)\).
The symmetric bipartite form of \(Q\) thus satisfies
Lemma~\ref{lem:scaling}, since
\(\varepsilon(\log N)^2\to0\).
It gives positive row and column multipliers \(u,v\) with
\[
 \widetilde D=\diag(u)Q\diag(v),\quad
 \max(|u_i-1|,|v_j-1|)
                     \leq C_\delta\varepsilon\log N.
\]

The total scaling cost requires compensation, not just
the maximum multiplier. In the symmetric bipartite form
write \(z=(u-\one,v-\one)\), and denote its row degrees by \(d_i\).
Sum the exact scaling equations
\(Q_{\rm bip}z+d\mathbin{\odot}z+
 z\mathbin{\odot}(Q_{\rm bip}z)=\one-d\)
over all \(2n\) coordinates:
\[
 \sum_i d_i z_i
       =r-s_{I,J}-\tfrac12 z^\top Q_{\rm bip}z .
\]
Using \(1-d_i\leq\varepsilon\), the maximum bound on \(z\),
and the full quadratic term, we obtain
\[
 \sum_i\log u_i+\sum_j\log v_j
 =r-s_{I,J}
        +O_\delta(r^2A^2(\log N)^2/N).
\]
This retains the overlap mass \(s_{I,J}\) and the nonlinear
compensation term.

Put \(\widetilde B=\widetilde D-J_n/n\).
After padding its row and column sets by zero, its Frobenius
distance from \(B\) is at most
\[
 C_\delta\left(\sqrt{rA/N}+r/N+
                     rA^{3/2}\log N/N\right).
\]
The first term removes the actual rows and columns, the
second changes the constant block, and the third includes
both scaling multipliers.
Since \(\|B\|_F\leq\sqrt A\), the trace-norm distance of the
squared singular-value matrices is at most
\[
 C_\delta\left(A\sqrt{r/N}+r\sqrt A/N+
                                      rA^2\log N/N\right).
\]
Both squared matrices have norm bounded away from one.
The trace-norm log-determinant inequality therefore shows
that their determinant factors differ by
\(\exp(O_{\delta,K}(N^{-1/4}))\).
All displayed errors are powers of \(\log N\) times \(N^{-1/2}\)
or \(N^{-1}\), uniformly over the supports.

Apply Theorem~\ref{thm:growingpermanent} to \(D\) and
\(\widetilde D\). The latter has a fixed gap and entry bound
at most \(3A/n\), fitting a possibly larger fixed
polylogarithmic exponent. Exactly,
\[
 \operatorname{per}Q=
                 \operatorname{per}\widetilde D/(\prod_i u_i\prod_jv_j).
\]
Stirling's formula gives
\[
 \log\frac{n!/n^n}{N!/N^N}=r+O(r^2/N).
\]
Its leading \(r\) cancels the full leading scaling cost.
The overlap term is \(O(r^2A/N)\); all the remaining scaling,
determinant and factorial errors are \(o(N^{-1/8})\).
The two relative permanent errors are \(O(N^{-1/8})\).
This proves the stated uniform \(O(N^{-1/9})\) comparison.
\end{proof}
